\chapter{The Navier--Stokes release: exact scope and proof interfaces}
\label{ch:navier-stokes-release}

\reviewstatus{The 166-page released manuscript has been read in full across the
two mathematical review passes recorded in the source ledger. The repository
revision, final theorem statements, dependency graph, and selected implementation
modules have been inspected. Reading a mathematical proof, rebuilding a Lean
target, and independently checking an exported kernel certificate are distinct
operations; their actual outcomes are recorded separately. The propositions
proved in this chapter explain mathematical interfaces of the construction.
They do not replace the analytic construction imported from the release.}

\section{What the released theorem says}

OpenAI's release of 8 September 2026 concerns a smooth, externally forced,
three-dimensional incompressible flow. It supplies a manuscript and Lean
formalization. The Clay Mathematics Institute's response of 11 September
acknowledges the apparent resolution while describing an ongoing evaluation
process. These facts must be distinguished from an award decision
\citep{NSOpenAIAnnouncement2026,NSClayAnnouncement2026}.

Let $\nu>0$ and define the momentum residual
\begin{equation}
 \mathcal R_\nu(u,p)
 =\partial_tu+(u\cdot\nabla)u-\nu\Delta u+\nabla p.
 \label{eq:ns-residual}
\end{equation}
The source's main conclusion can be recorded without changing its quantifiers:
\begin{align}
 &f\in C_c^\infty(\R^3\times(0,\infty);\R^3),\qquad
 u,p\in C^\infty(\R^3\times[0,1)),\nonumber\\
 &\mathcal R_\nu(u,p)=f,\quad \nabla\cdot u=0,\quad u(\cdot,0)=0,
 \label{eq:ns-release-equation}\\
 &\supp u(\cdot,t)\cup\supp p(\cdot,t)\subset K
 \quad(0\leq t<1),\nonumber\\
 &\sup_{t<1}\|u(t)\|_2<\infty,\qquad
 \limsup_{t\uparrow1}\|u(t)\|_\infty=\infty.
 \label{eq:ns-release-growth}
\end{align}
Here $K$ is compact. With the same datum and force there is no global smooth
whole-space solution with uniformly bounded kinetic energy. A corresponding
periodic result is also established \citep[Theorem 1.1 and Corollary 10.6]{NSOpenAI2026}.

In Fefferman's formulation, alternatives (A) and (B) impose $f=0$, whereas
(C) and (D) permit a smooth force satisfying the specified decay conditions.
The released statements address (C) and (D). They do not logically negate (A)
or (B): an existential assertion with $f\neq0$ does not supply an example with
$f=0$. Pressure periodicity is part of the periodic formulation, including its
erratum \citep{NSFefferman}. The separate Euler result concerns a different
equation, obtained by setting viscosity to zero, and cannot be transferred to
positive viscosity by deleting the viscous estimates.

This distinction changes the appropriate goal of the present extension. A
more informative presentation can expose the cancellation, support, analytic
norms, and formal dependency interfaces. A stronger theorem would require a
new mathematical conclusion with all its additional hypotheses discharged.
The material below provides the former and identifies precisely where the
source's construction supplies the latter's indispensable analytic inputs.

\section{An exact algebra for residual correction}

The same distinction between a complex and an exact complex encountered
earlier reappears here. A differential identity constrains what a correction
can do. It does not prove that a correction of the necessary size, support,
and regularity exists.

\begin{proposition}[Quadratic residual identity]
\label{prop:ns-residual-increment}
For smooth $u,w,p,\pi$, with $\nabla\cdot w=0$,
\begin{align*}
 \mathcal R_\nu(u+w,p+\pi)
 &=\mathcal R_\nu(u,p)+L_{u,\nu}(w,\pi)
       +\nabla\cdot(w\otimes w),\\
 L_{u,\nu}(w,\pi)
 &=\partial_tw+(u\cdot\nabla)w+(w\cdot\nabla)u
       -\nu\Delta w+\nabla\pi .
\end{align*}
The tensor convention is $(w\otimes w)_{ij}=w_iw_j$ and
$(\nabla\cdot S)_i=\sum_j\partial_jS_{ij}$.
\end{proposition}
\begin{proof}
All terms other than advection are linear. Expanding the latter gives
$(u+w)_j\partial_j(u_i+w_i)$, with the four products obtained by choosing
one of $u,w$ in each factor. Moreover,
$\sum_j\partial_j(w_iw_j)=\sum_jw_j\partial_jw_i+
w_i\sum_j\partial_jw_j=(w\cdot\nabla)w_i$. This proves every component
of the identity. No averaging, limiting argument, or neglected term occurs.
\end{proof}

In particular, solving $L_{u,\nu}(w,\pi)=-\mathcal R_\nu(u,p)$ leaves the
quadratic residual $\nabla\cdot(w\otimes w)$. A numerical decrease in the
linear residual alone is therefore an incomplete verification. Conversely,
if a background satisfies
$\mathcal R_\nu(u,p)=-\nabla\cdot S+E$, a fluctuation with averaged
covariance $S$ can cancel the selected mean stress while creating other
linear and oscillatory terms. Each of those terms needs a place in the next
residual estimate.

\begin{proposition}[A transverse wave and its covariance]
\label{prop:ns-transverse-wave}
Let $a,k\in\R^3$ with $a\cdot k=0$ and put
$w(x,\phi)=a\cos(k\cdot x+\phi)$. Then
\[
 \nabla\cdot w=0,\qquad (w\cdot\nabla)w=0,\qquad
 \frac1{2\pi}\int_0^{2\pi}w\otimes w\,d\phi
       =\frac12a\otimes a.
\]
The vanishing self-advection statement requires constant $a$ and $k$.
\end{proposition}
\begin{proof}
Differentiating the cosine gives
$\nabla\cdot w=-(a\cdot k)\sin(k\cdot x+\phi)$ and
$(w\cdot\nabla)w=-(a\cdot k)a\cos(k\cdot x+\phi)
\sin(k\cdot x+\phi)$. Both vanish. Integration of
$\cos^2s=(1+\cos2s)/2$ proves the tensor identity.
If the amplitude or phase gradient varies, the product rule introduces
additional derivatives, so this proof gives no cancellation for them.
\end{proof}

The covariance is an element of $\Sym^2(\R^3)$, with nonnegative quadratic
form $v\mapsto (a\cdot v)^2/2$. The positivity is a concrete use of a
symmetric tensor power. A rank-one tensor representation does not by itself
solve the vector-valued stress-divergence equation; the allowed polarization
directions and the differential constraints still matter. The exterior-power
test for independence of two stress directions is developed in
Chapter~\ref{ch:fluid-structures}.

\begin{proposition}[Averaging and nonlinear closure]
\label{prop:ns-average-defect}
Let $P$ be a linear averaging operator on a function algebra with $P1=1$.
Assume $P^2=P$ and the conditional-expectation identity
$P(ag)=aP(g)$ whenever $Pa=a$. If $v=Pu$ and $w=u-v$, then
\[
 P(u\otimes u)=v\otimes v+P(w\otimes w).
\]
If $P$ is positive, the final term is positive semidefinite pointwise in
the averaged variables.
\end{proposition}
\begin{proof}
Idempotence gives $Pw=0$. Expand $(v+w)\otimes(v+w)$. The two mixed
terms have zero average by the assumed multiplication identity, and
$P(v_i v_j)=v_iv_j$. For a constant test vector $b$,
$b^TP(w\otimes w)b=P((b\cdot w)^2)\geq0$.
\end{proof}

An orthogonal projection onto an arbitrary collection of modes need not have
the multiplication property in this proposition. Likewise, averaging an
extended function $F(x,Y)$ over independent $Y$ and then substituting a phase
$Y=\Phi(x)$ is different from first restricting and averaging the resulting
physical function. If $F(x,Y)=\cos(2\pi(Y-x))$ on one-dimensional tori,
then $\int F(x,Y)dY=0$, whereas $F(x,x)=1$. This elementary example prevents
an unjustified replacement of a pointwise physical residual by an auxiliary
mean. The phase derivatives and support geometry in a complete construction
are precisely what allow these two kinds of expression to be related.

\section{Concentration, energy, and the information lost by norms}

\begin{proposition}[Energy estimate from rest]
\label{prop:ns-force-energy}
Let $u,p$ be a smooth solution on $[0,T]$ of
$\mathcal R_\nu(u,p)=f$, $\nabla\cdot u=0$, $u(0)=0$.
Assume periodicity, or sufficient decay for the integrations by parts
below; fixed compact spatial support of $u,p$ is sufficient. If
$F(t)=\int_0^t\|f(s)\|_2\,ds$, then
\[
 \|u(t)\|_2^2+2\nu\int_0^t\|\nabla u(s)\|_2^2\,ds\leq F(t)^2.
\]
\end{proposition}
\begin{proof}
The pressure pairing is
$\int u\cdot\nabla p=-\int p\nabla\cdot u=0$.
Advection contributes
$\int u\cdot(u\cdot\nabla)u=
\frac12\int\nabla\cdot(u|u|^2)=0$. Thus
\[
 \frac12\frac d{dt}\|u\|_2^2+\nu\|\nabla u\|_2^2=\langle f,u\rangle.
\]
Write $y_\delta=(\|u\|_2^2+\delta^2)^{1/2}$, $\delta>0$.
Then $y_\delta'\leq\|f\|_2$, so
$y_\delta(t)\leq\delta+F(t)$. Passing to $\delta\downarrow0$ gives
$\|u(t)\|_2\leq F(t)$. Integrating the energy identity once more,
\[
 \|u(t)\|_2^2+2\nu\int_0^t\|\nabla u\|_2^2
 \leq2\int_0^t F(s)F'(s)\,ds=F(t)^2.
\]
The argument on each $T<1$, followed by monotone convergence of the
nonnegative dissipation integral, also applies up to a singular time.
\end{proof}

This estimate controls an integral, not the maximum of the integrand.
Information reduction becomes mathematically consequential at this exact
point: a small spatial set can carry a large amplitude while making a small
contribution to $L^2$ energy.

\begin{proposition}[Anisotropic concentration test]
\label{prop:ns-anisotropic-scaling}
Let $0<h<1/6$, $\tau>0$, and $W\in C_c^\infty(\R^3;\R^m)$ be nonzero.
Define the comparison family
\[
 V_\tau(x_1,x_2,x_3)=\tau^{-1/2-h}
 W\left(\frac{x_1}{\tau^{1/2}},
        \frac{x_2}{\tau^{1/2}},
        \frac{x_3}{\tau^{1/2-h}}\right).
\]
For $1\leq p<\infty$,
\[
 \|V_\tau\|_p^p
  =\tau^{3/2-h-p(1/2+h)}\|W\|_p^p,
 \qquad
 \|V_\tau\|_\infty=\tau^{-1/2-h}\|W\|_\infty.
\]
Its energy tends to zero, while its maximal amplitude diverges. Moreover,
$\int_0^{\tau_0}\|\nabla V_\tau\|_2^2\,d\tau<\infty$ for every finite
$\tau_0>0$. This comparison family is not asserted to solve Navier--Stokes.
\end{proposition}
\begin{proof}
The change of variables has determinant $\tau^{3/2-h}$. This proves the
norm identities. At $p=2$ the exponent is $1/2-3h>0$. Each of the first
two spatial derivatives contributes $\tau^{-1/2}$, so its squared norm has
exponent $-1/2-3h>-1$. The axial derivative contributes
$\tau^{-1/2+h}$ and has squared-norm exponent $-1/2-h>-1$.
Both powers are integrable at zero. The supremum formula follows because
the coordinate dilation is onto.
\end{proof}

For this family the critical spatial integrability exponent is
\[
 p_*(h)=\frac{3/2-h}{1/2+h}<3.
\]
At $p=3$, $\|V_\tau\|_3^3=\tau^{-4h}\|W\|_3^3$, which diverges.
Thus the scaling test does not contradict a regularity theorem requiring
a bounded critical $L^3$ norm. Such a theorem is a substantially stronger
statement than an energy bound; the endpoint theorem for the unforced
Cauchy problem is due to Escauriaza, Seregin, and \v Sver\'ak
\citep{NSEscauriazaSereginSverak2003}. Our comparison family proves only
the logical insufficiency of the weaker norm information, independently
of any existence construction.

\section{Similarity coordinates as an exact change of variables}

The geometry relevant to the release can be organized using the positive
variables
\[
 \tau=1-t,\qquad D=\tfrac12-h,\qquad
 \tau=q(1-\eta^2),\quad z=q^D\eta,\quad X=\frac{r^2}{2q},
 \quad 0<h<\tfrac12.
\]
The next calculation verifies the differential interface rather than
assuming that symbolic substitution preserves all derivatives.

\begin{proposition}[Coordinate existence and derivative weights]
\label{prop:ns-similarity-chain-rule}
For $\tau>0$ and $z\in\R$ there is a unique
$q>|z|^{1/D}$ with $q-z^2q^{2h}=\tau$. Set $\eta=zq^{-D}$,
$d=1-\eta^2$, and $L=1-2h\eta^2$. Then $|\eta|<1$, $L\geq1-2h>0$,
and for a smooth profile $G(X,\eta)$ and any real $b$,
\begin{align*}
 \partial_t(q^bG)&=q^{b-1}L^{-1}
     (-bG+D\eta\partial_\eta G+X\partial_XG),\\
 \partial_z(q^bG)&=q^{b-D}L^{-1}
     (2b\eta G+d\partial_\eta G-2\eta X\partial_XG).
\end{align*}
\end{proposition}
\begin{proof}
On $q>|z|^{1/D}$, the function $q-z^2q^{2h}$ begins at zero,
tends to infinity, and has derivative
$1-2hz^2q^{2h-1}\geq1-2h>0$. The intermediate-value theorem and strict
monotonicity give existence and uniqueness. The implicit-function theorem
gives smoothness, and differentiation yields
\[
 q_t=-L^{-1},\quad q_z=2\eta q^{1-D}/L,\quad
 \eta_t=D\eta/(qL),\quad\eta_z=d/(q^DL),
\]
\[
 X_t=X/(qL),\qquad X_z=-2\eta X/(q^DL).
\]
The two formulas follow by the product and chain rules. The endpoint $z=0$
also satisfies them, because the implicit derivative remains nonzero.
\end{proof}

On bounded profile ranges, a transverse derivative has weight $q^{-1/2}$,
an axial derivative has weight $q^{-D}$, and a time derivative has weight
$q^{-1}$. The denominators remain bounded because $L\geq1-2h$.
Constants can depend on the derivative order and fixed profile parameters.
They cannot be allowed to hide an extra loss growing with the approximation
stage if one later intends to sum an infinite correction sequence.

The apparent cylindrical singularity at $r=0$ must also be resolved.
For instance, if the azimuthal velocity is $r\,a(r^2,z,t)$, then
\[
 r\,a(r^2,z,t)e_\theta
   =a(x_1^2+x_2^2,z,t)(-x_2,x_1,0)
\]
is smooth in Cartesian coordinates. Boundedness of a cylindrical component
alone would not establish this. This is a local regularity condition on a
representation, comparable to requiring that coordinate expressions of a
section glue smoothly; it is not supplied by the topology of the domain.

\section{A quadratic relation that controls analytic derivative loss}

An especially useful connection to the project's algebraic methods occurs
in the analytic linear recursion. The relevant feature is a forbidden
composition of derivative-carrying blocks. We isolate its mechanism in a
statement independent of fluid terminology.

\begin{lemma}[Words with no adjacent derivative letters]
\label{lem:ns-forbidden-words}
A word of length $k$ in the letters $N,D$ with no consecutive $D$ has at
most $\lceil k/2\rceil$ occurrences of $D$. The number with exactly $j$
occurrences is $\binom{k-j+1}{j}$, interpreted as zero when $j>k-j+1$.
The total number $a_k$ satisfies $a_0=1$, $a_1=2$, and
$a_k=a_{k-1}+a_{k-2}$ for $k\geq2$.
\end{lemma}
\begin{proof}
Between $j$ derivative letters there are at least $j-1$ other letters, so
$2j-1\leq k$. If their positions are
$1\leq i_1<\cdots<i_j\leq k$ with $i_{r+1}\geq i_r+2$, the shifted
positions $i_r-(r-1)$ form an ordinary $j$-element subset of
$\{1,\ldots,k-j+1\}$. This is a bijection, proving the binomial formula.
For the recurrence, an admissible word ending in $N$ is an arbitrary
admissible word of length $k-1$ followed by $N$. One ending in $D$, when
$k\geq2$, ends in $ND$ and corresponds to a word of length $k-2$.
These two classes are disjoint and exhaustive.
\end{proof}

\begin{proposition}[Sparse Volterra systems]
\label{prop:ns-sparse-volterra}
Let $E=E_0\oplus E_1$ be a finite-dimensional normed complex vector space.
Fix a compact real interval $I$ and let
$S_\rho=\{z\in\C:\dist(z,I)<\rho\}$. Consider an integral equation
on $0\leq s\leq a$ and this complex parameter neighborhood,
\[
 W=Gf+G(A_0W+A_1\partial_\eta W),\qquad
 (Gg)(s)=\int_0^s K(s,r)g(r)\,dr.
\]
Assume $\|K(s,r)\|\leq1$, $K$ preserves the two summands and is independent
of $\eta$; the coefficients and kernel are continuous in their radial
variables; $A_0,A_1,f$ are uniformly bounded and holomorphic in $\eta$;
and every value and every $\eta$-derivative of $A_1$ maps $E_0$ into $E_1$
and annihilates $E_1$. On each smaller neighborhood $S_{\rho'}$ with
$0<\rho'<\rho$, the Picard series converges for every finite $a$.
It gives the unique solution in the class of solutions bounded and
holomorphic on a parameter neighborhood strictly larger than the one
on which they are compared.
\end{proposition}
\begin{proof}
Expand the $k$-fold iterate of $G(A_0+A_1\partial_\eta)$ into words.
The kernel between adjacent coefficients preserves $E_0,E_1$. Therefore
\[
 A_1(s)K(s,r)A_1(r)=0,
 \qquad A_1(s)K(s,r)\partial_\eta A_1(r)=0.
\]
The second identity is necessary: in a composition of two derivative
operators the first derivative can strike the second coefficient. It
vanishes for the same block reason. Thus a nonzero word cannot contain
two adjacent derivative operators, even after expanding the product rule.
By Lemma~\ref{lem:ns-forbidden-words}, at most
$p_k=\lceil k/2\rceil$ Cauchy radius losses are required.

Put $\Delta=\rho-\rho'>0$. Distribute that loss in $p_k$ equal parts.
The Cauchy estimate for one parameter derivative costs at most
$p_k/\Delta$; when $p_k=0$ use cost one. Operator norms absorb the
coefficient bounds, and the nested integration domain has volume
$a^{k+1}/(k+1)!$. Bounding the number of words by $2^k$ gives, after
enlarging a constant $C$,
\[
 \bigl\|[G(A_0+A_1\partial_\eta)]^kGf\bigr\|_{S_{\rho'}}
 \leq\frac{C^{k+1}a^{k+1}}{(k+1)!}
       \max\{1,p_k/\Delta\}^{p_k}.
\]
Stirling's bound shows that the $k$th root of the right-hand side is
$O(k^{-1/2})$, tending to zero. Hence the series converges uniformly on
the smaller neighborhood. Applying the same estimate after one additional
radius restriction permits differentiation in $\eta$ and termwise
substitution into the integral equation. If $W_1-W_2=G(A_0+A_1\partial_\eta)
(W_1-W_2)$, iteration and the same estimate make the difference zero on
the smaller neighborhood. Analytic continuation gives the stated uniqueness.
\end{proof}

The source's six-component linear recursion has the splitting
$E_0=\C^4$, $E_1=\C^2$, with diagonal radial kernels. Its use of the
above relation is traceable to Lemma 5.1, rather than inferred from a
superficial resemblance to a quiver \citep{NSOpenAI2026}. The relation
has three precise interpretations. It eliminates forbidden operator words;
it limits derivative loss; and, with the integration simplex, it makes a
convergent analytic solve possible on a common radial interval.

The language of admissible words is regular: a finite automaton needs only
to remember whether the preceding letter was $D$, plus a rejecting state.
The analytic conclusion nevertheless requires unbounded iteration and
quantitative estimates. A regular syntax does not make the represented
analytic task finite-dimensional.

Nor does this calculation establish a gentle algebra. The abstract
one-vertex algebra $k\langle N,D\rangle/(D^2)$ is infinite-dimensional,
and the arrow $N$ has two allowed successors. It therefore fails a usual
gentle-algebra continuation condition. The honest bridge is to a
representation with a quadratic zero relation and a controlled path
language. Any claim of a particular gentle presentation would require
a separate equivalence or construction.

\section{From formal orders to actual smooth fields}

The order of a term in a formal expansion must not be confused with a
convergence estimate for the coefficients. Suppose coefficients of order
$j$ have constants growing faster than every exponential. Positive powers
of a small parameter do not automatically make their unrestricted sum
converge. A suitable realization uses cutoffs chosen after the coefficients
are known.

\begin{lemma}[A diagonal realization estimate]
\label{lem:ns-diagonal-realization}
Let $0<q<q_0\leq1$ on an open domain, with $q$ smooth and bounded below
on every compact subset. Suppose smooth increments $Z_j$, $j\geq1$,
satisfy
\[
 |Z_j|_m\leq C_{j,m}q^{g_j-\ell_m}(1+|\log q|)^{P_{j,m}},
 \qquad 0<g_1\leq g_2\leq\cdots\longrightarrow\infty,
\]
where $|\cdot|_m$ is the pointwise maximum of coordinate derivatives through order
$m$, and $\ell_m$ is independent of $j$. Assume that differentiating a
cutoff $\chi(aq)$ costs a power of $q^{-1}$ depending only on derivative
order, uniformly in $a$ on its transition region. Then one can choose
$a_{j+1}\geq2a_j$ so that the cutoff sum
$Z=\sum_{j\geq1}\chi(a_jq)Z_j$ is locally finite and smooth, and its
tail after a fixed initial block satisfies
\[
 \left|Z-\sum_{j=1}^JZ_j\right|_m
 \leq2^{-J}q^{g_{J+1}/2-L_m}
 \quad\left(J\geq\max\{1,m\},\ q<\frac1{2a_J}\right).
\]
Here $\chi=1$ on $[0,1/2]$, $\chi=0$ on $[1,\infty)$, and $L_m$ is
independent of $J$.
\end{lemma}
\begin{proof}
The product rule and the cutoff hypothesis bound the $m$th seminorm of
the $j$th summand by
$\widehat C_{j,m}q^{g_j-L_m}(1+|\log q|)^{\widehat P_{j,m}}$.
At stage $j$ impose this finite list of requirements for $0\leq m\leq j$:
\[
 \widehat C_{j,m}q^{g_j/2}(1+|\log q|)^{\widehat P_{j,m}}
 \leq2^{-j}\quad(0<q\leq a_j^{-1}).
\]
It is possible because every positive power of $q$ dominates every fixed
power of $|\log q|$ at zero. Increase $a_j$ further to preserve its geometric
growth. On a compact subset where $q\geq\delta>0$, terms with
$a_j>\delta^{-1}$ vanish, proving local finiteness and smoothness.
If $q<(2a_J)^{-1}$, the first $J$ cutoffs equal one. For $j>J\geq m$,
the bound just imposed and $g_j\geq g_{J+1}$ give
\[
 \sum_{j>J}2^{-j}q^{g_j/2-L_m}
 \leq2^{-J}q^{g_{J+1}/2-L_m},
\]
as required.
\end{proof}

To preserve a differential constraint, apply the cutoffs to suitable
potentials. In three-dimensional Euclidean space,
\[
 \nabla\times(\chi(a_jq)A_j)
 =\chi(a_jq)\nabla\times A_j
   +a_j\chi'(a_jq)\nabla q\times A_j.
\]
The second term is a real part of the field and must be retained. Taking
divergence of the full curl yields zero. Cutting off an already
divergence-free velocity directly would instead give
$\nabla\cdot(\chi u)=\nabla\chi\cdot u$, which need not vanish.

\begin{proposition}[All-order residual realization]
\label{prop:ns-residual-flatness}
Let $U^{[J]}$ be finite approximations and $U$ their cutoff realization.
Assume their derivatives have polynomial growth in $q^{-1}$ with
exponent losses independent of $J$, and that the tail estimate of
Lemma~\ref{lem:ns-diagonal-realization} applies through every fixed order.
Let $F$ be a fixed differential polynomial of finite order and degree whose
coefficient derivatives have polynomial growth in $q^{-1}$. If
\[
 |F(U^{[J]})|_m
 \leq C_{J,m}q^{\rho_J-K_m}(1+|\log q|)^{P_{J,m}}+E_{J,m},
 \quad \rho_J\to\infty,
\]
with $K_m$ independent of $J$ and each fixed $E_{J,m}=O(q^N)$ for every
$N$, then $|F(U)|_m=O(q^N)$ for every $m,N$.
\end{proposition}
\begin{proof}
Fix $m,N$. Let $s,d$ bound the differential order and polynomial degree
of $F$. For products use the exact identity
\[
 \prod_{i=1}^d a_i-\prod_{i=1}^d b_i
 =\sum_{r=1}^d
  \left(\prod_{i<r}a_i\right)(a_r-b_r)
  \left(\prod_{i>r}b_i\right).
\]
The product rule then bounds $F(U)-F(U^{[J]})$ in seminorm $m$ by a
fixed power of $q^{-1}$ times the tail seminorm $m+s$ and a polynomial
in the two field seminorms $m+s$. All exponent losses depend on $m,s,d$,
not $J$. Choose $J$ large enough that both $\rho_J$ and $g_{J+1}/2$
exceed $N$ plus these losses and one extra unit. Keep that $J$ fixed.
Shrink the $q$-neighborhood so all its finitely many constants and logarithms
are absorbed by that extra unit, and the first $J$ cutoffs equal one.
The finite residual, its fixed flat remainder, and the residual difference
are then all $O(q^N)$.
\end{proof}

The order of these quantifiers is essential:
\[
 \forall(m,N)\ \exists J\ \exists\delta,C\ 
  \forall q\in(0,\delta):\ |F(U)|_m\leq Cq^N.
\]
It does not assert one $J$ that controls every derivative. It does not
require summing the flat remainders $E_{J,m}$ over $J$. It also does not
assert convergence of the uncut formal series. This distinction is the
analytic counterpart of keeping a filtered algebra, its associated graded
object, and a completed realization separate.

\section{The finite correction cycle as a verifiable dependency graph}

The source's construction provides the hypotheses needed above through
separate analytic modules. The following locator table records where to
check them; the table is a dependency guide, not a substitute proof
\citep{NSOpenAI2026,NSLean2026}.

\begin{longtable}{p{.20\textwidth}p{.40\textwidth}p{.30\textwidth}}
\toprule
Interface & Mathematical obligation & Source location\\
\midrule\endhead
Background & Smooth axis; exact exterior; supported compatible stress;
positive cone margins & Theorem 4.6; Appendices A--C\\
Order recursion & A common radial interval, moment repair, derivative losses
independent of formal order & Lemmas 5.1--5.4; Proposition 5.5\\
Phase geometry & Support separation and compatible finite covering coordinates
& Section 6\\
Oscillatory solve & Transverse evolution, damping, real covariance, full curl
correction & Propositions 7.2, 7.5--7.6; Lemma 7.7\\
Mean solve & Compact primitives, auxiliary inverse, five moment equations
& Section 8\\
Iteration & Full residual recomputation, fixed-domain gains, all physical
derivatives & Propositions 9.3--9.9\\
Completion & Smooth forcing, energy, comparison, all viscosities, periodization
& Section 10\\
\bottomrule
\end{longtable}

One can encode a correction state as a typed record rather than a loose
list of estimated errors. Its fields should include a velocity, pressure,
potential representation, support certificates, exact preserved moments,
residual decomposition, and derivative bounds. A correction operation
must produce a new such record. The pressure and residual are recomputed
from the updated velocity, so one cannot reuse a previous covariance after
altering a wave amplitude. This record-based reading is directly compatible
with the typed graph approach developed earlier in the monograph.

The numerical exponents also permit an elementary exact check. Put
$\kappa=10^{-5}$, $\sigma_j=1/5+j/10$,
$B_j=1/2+\sigma_j$, and $C_j=1+\sigma_j$.
All asserted gains of the source's cycle must exceed $1/10$ after
derivative losses. For example,
\[
 \frac{18}{100}-2\kappa>\frac{17}{100}>\frac1{10},\qquad
 \sigma_j-3\kappa\geq\frac15-3\kappa>\frac{17}{100},
\]
\[
 \frac12-2\kappa>\frac1{10},\qquad
 \frac9{10}-4\kappa>\frac1{10}.
\]
These rational inequalities certify arithmetic slack once the functional
estimates are supplied. They do not establish those functional estimates.
This is a useful separation for learned proof search: a model may propose
an ordering or an exponent assignment, while exact arithmetic and analytic
lemmas discharge different obligations.

\section{Localization and smooth force extension}

\begin{proposition}[Localization without a divergence defect]
\label{prop:ns-potential-localization}
Suppose $A$, $B e_\theta$ and $p$ are smooth in Cartesian coordinates,
put $u=\nabla\times A+B e_\theta$, and assume $B$ is independent of $\theta$.
Let $c$ be a smooth axisymmetric scalar.
Set
\[
 \widetilde u=\nabla\times(cA)+cBe_\theta,
 \qquad \widetilde p=cp,\qquad a=\nabla c\times A.
\]
Then $\widetilde u=cu+a$ and $\nabla\cdot\widetilde u=0$.
Its residual is exactly
\begin{align*}
 \mathcal R_\nu(\widetilde u,cp)
 ={}&c\mathcal R_\nu(u,p)
   +(\partial_tc-\nu\Delta c)u-2\nu\sum_j(\partial_jc)\partial_ju
       +p\nabla c\\
 &+(c^2-c)(u\cdot\nabla)u+c(u\cdot\nabla c)u\\
 &+\partial_ta-\nu\Delta a+c(u\cdot\nabla)a
      +(a\cdot\nabla)(cu)+(a\cdot\nabla)a.
\end{align*}
\end{proposition}
\begin{proof}
The curl product rule gives $\widetilde u=cu+a$. Divergence of a curl
is zero; the other term has divergence $r^{-1}\partial_\theta(cB)=0$
away from the axis and hence everywhere by Cartesian smoothness.
Expand $\partial_t(cu)$, $\Delta(cu)$, and $\nabla(cp)$ by the product
rule. Finally,
$(cu\cdot\nabla)(cu)=c^2(u\cdot\nabla)u+c(u\cdot\nabla c)u$.
Adding the three advection terms containing $a$ yields the formula.
\end{proof}

The displayed expression explains why a spatial cutoff requires endpoint
regularity of the actual potentials as well as the velocity. On a region
where $c=1$ the extra terms vanish. On its transition region every term
involving the uncut fields must remain smooth through the terminal time.
A proof that estimates only $c\mathcal R_\nu(u,p)$ has omitted real terms.

\begin{lemma}[Realization of terminal force jets]
\label{lem:ns-terminal-jets}
Let $f(x,t)$ be smooth for $t<T$, supported in a fixed compact set
$K\subset\R^d$. Suppose every $\partial_x^\alpha\partial_t^jf$
has a uniform limit $\partial_x^\alpha F_j(x)$ as $t\uparrow T$,
where $F_j\in C_c^\infty(\R^d)$ are supported in $K$.
Then $f$ has a smooth extension to $t\geq T$ supported in
$K\times[T,T+1]$ on that side. If it already vanishes near an initial
time, the combined extension retains that property.
\end{lemma}
\begin{proof}
Choose a smooth cutoff $\chi$, one near zero and zero for arguments
at least one, and for $s=t-T\geq0$ define
\[
 f_+(x,T+s)=\sum_{j=0}^\infty\chi(b_js)\frac{s^j}{j!}F_j(x).
\]
For fixed $\alpha,m$, the product rule and $s\leq b_j^{-1}$ on the
summand's support give a bound
\[
 C_{j,m}b_j^{m-j}\|F_j\|_{C^{|\alpha|}}.
\]
Fix $b_0=1$. For $j\geq1$, choose increasing $b_j\geq1$ so that this is at most $2^{-j}$ whenever
$|\alpha|+m\leq\lfloor j/2\rfloor$. Each step imposes finitely many
constraints with $m-j<0$, so the choice exists. Every differentiated tail
converges uniformly. At $s=0$, the $m$th time derivative of a fixed
summand is $F_m$ when $j=m$ and zero otherwise. Uniform differentiated
convergence gives the desired jets.

Uniform convergence of the original mixed derivatives makes their limits
compatible: pass to the limit in the fundamental theorem of calculus
along a coordinate segment, and then along a time segment. The two
pieces therefore glue smoothly at $T$. All summands retain support in
$K$, and vanish for $s\geq1$; uniform convergence of every derivative
also proves smoothness at that support boundary.
\end{proof}

Flatness at a single singular point is only part of this lemma's hypotheses.
One must additionally control all terminal derivatives away from that point,
including the cutoff transition region. Conversely, vanishing of every jet
does not force a smooth function to vanish near the point: the function
$e^{-1/s^2}$ for $s>0$, extended by zero for $s\leq0$, is the standard
counterexample. Substituting ``analytic'' for ``smooth'' would change the
problem and invalidate this extension method.

\section{Comparison, quantitative residual checks, and proof strength}

An explicit singular solution rules out a global competitor only after
uniqueness has been established in the competitor's class. The whole-space
class in the Clay problem permits no arbitrary decay assumption on the
pressure. That is why a periodic energy calculation is not by itself a
complete proof of alternative (C). The source's Lemma 10.5 handles this
point using a pressure-gradient argument and expanding spatial cutoffs;
its hypotheses must be preserved in any proof adapter
\citep{NSOpenAI2026}.

The following separate estimate is useful for the project's computable
fluid examples and for any learned approximate field.

\begin{proposition}[Residual-to-solution estimate on the flat torus]
\label{prop:ns-residual-comparison}
Let $u,v$ be smooth divergence-free periodic vector fields on $[0,T]$,
and let periodic $p,q$ satisfy
\[
 \mathcal R_\nu(u,p)=f,\qquad
 \mathcal R_\nu(v,q)=f+r.
\]
Write $w=v-u$ and $K(t)=\|\nabla u(t)\|_{L^\infty,\mathrm{op}}$.
If $K$ is integrable, then
\[
 \|w(t)\|_2\leq
 e^{\int_0^tK(a)\,da}\|w(0)\|_2
 +\int_0^t e^{\int_s^tK(a)\,da}\|r(s)\|_2\,ds.
\]
In particular, equal data and $r=0$ imply $v=u$.
\end{proposition}
\begin{proof}
Subtract the equations in the form
\[
 \partial_tw+(v\cdot\nabla)w+(w\cdot\nabla)u
 -\nu\Delta w+\nabla(q-p)=r.
\]
Periodic integration and zero divergence remove the transport and pressure
pairings. Thus
\[
 \tfrac12(\|w\|_2^2)'+\nu\|\nabla w\|_2^2
 \leq K\|w\|_2^2+\|r\|_2\|w\|_2.
\]
With $y_\delta=(\|w\|_2^2+\delta^2)^{1/2}$ one has
$y_\delta'\leq Ky_\delta+\|r\|_2$. Multiplication by
$e^{-\int_0^tK}$ and integration prove the bound after
$\delta\downarrow0$.
\end{proof}

This inequality gives a precise role for a quantum or classical learner:
it may produce $v$, but a certificate needs its divergence defect, momentum
residual, initial error, and an appropriate stability coefficient. The
estimate does not remain effective near blowup merely because the sampled
residual is small. A pointwise grid evaluation does not bound the continuum
$L^2$ residual without an interpolation or enclosure argument, and
$\int K$ can become large. Numerical fitting and a posteriori proof are
different stages of the same workflow.

The same issue explains why nearby classical results are not interchangeable.
Nonuniqueness in a weak class does not give blowup of a smooth solution;
a modified averaged nonlinearity is not the original Navier--Stokes
nonlinearity; and rough forcing is not smooth forcing. These distinctions
are visible respectively in the settings of Buckmaster--Vicol,
Tao, and Albritton--Bru\'e--Colombo
\citep{NSBuckmasterVicol2019,NSTao2016,NSAlbrittonBrueColombo2022}.

\section{Viscosity, periodization, and preservation of quantifiers}

\begin{proposition}[Viscosity conversion]
\label{prop:ns-viscosity-conversion}
If $\mathcal R_1(u,p)=f$ and $\nabla\cdot u=0$, then for $\nu>0$
\[
 u_\nu(x,t)=\sqrt\nu\,u(x/\sqrt\nu,t),\quad
 p_\nu(x,t)=\nu p(x/\sqrt\nu,t),\quad
 f_\nu(x,t)=\sqrt\nu\,f(x/\sqrt\nu,t)
\]
satisfy $\mathcal R_\nu(u_\nu,p_\nu)=f_\nu$ and zero divergence.
The singular time is unchanged. In three dimensions,
\[
 \|u_\nu(t)\|_2^2=\nu^{5/2}\|u(t)\|_2^2,
 \qquad
 \nu\|\nabla u_\nu(t)\|_2^2=\nu^{5/2}\|\nabla u(t)\|_2^2.
\]
\end{proposition}
\begin{proof}
Put $y=x/\sqrt\nu$. Each term of the momentum equation acquires a
factor $\sqrt\nu$: two spatial derivatives contribute $\nu^{-1}$,
one contributes $\nu^{-1/2}$, and the amplitudes contribute the stated
powers. The divergence equals $(\nabla_y\cdot u)(y,t)$. The Jacobian
of $x=\sqrt\nu y$ is $\nu^{3/2}$, while the velocity square contributes
$\nu$ and the gradient square contributes one. This proves both norm
identities. The time argument was not changed.
\end{proof}

\begin{proposition}[Disjoint-support periodization]
\label{prop:ns-periodization}
Suppose smooth whole-space fields $u,p,f$ solve Navier--Stokes with
viscosity $\nu$ and have common spatial support in a compact set
$K\Subset(-1/2,1/2)^3$. Then their sums over integer translates define
smooth periodic fields solving the same equation, with the same viscosity.
\end{proposition}
\begin{proof}
The sums are locally finite. The distinct translated compact supports have
positive separation, so at each point a neighborhood intersects at most
one support, unless the fields vanish there. Differential operators
commute with the sums. In the advection term every mixed product of
different translates vanishes, leaving the sum of the individual
advection terms. The same argument gives zero divergence and retains
periodicity of both velocity and pressure.
\end{proof}

For a larger compact support one first uses the parabolic scaling
\[
 \widetilde u(x,t)=\lambda u(\lambda x,\lambda^2(t-t_0)),\quad
 \widetilde p(x,t)=\lambda^2p(\lambda x,\lambda^2(t-t_0)),\quad
 \widetilde f(x,t)=\lambda^3f(\lambda x,\lambda^2(t-t_0)).
\]
Choosing $\lambda$ large puts the support inside the cube. If the original
fields vanish on an initial interval, they extend smoothly by zero before
$t_0$. For original singular time one, $t_0=1-\lambda^{-2}$ keeps the
new singular time equal to one. These conditions explain why periodizing
an arbitrary spatially decaying field would not suffice: overlapping copies
would create additional nonlinear interactions.

\section{Formal proof provenance and what this extension adds}

The pinned upstream revision is
\begin{center}\small
\texttt{f9e8bc5b38b6e212696e8a30e3e91517af887bbd}.
\end{center}
Its toolchain is Lean 4.34.0-rc2, with mathlib revision
\begin{center}\small
\texttt{85e3a25e006c35636f0e53b0e9296caca2685bc0}.
\end{center}
The project's earlier selected Lean declarations use a different toolchain.
They cannot be merged by silently substituting one compiler for the other.
An isolated verification script preserves the upstream pins.

The final upstream declarations are
\begin{quote}\small\ttfamily
NavierStokes.Comparator.navier\_stokes\_breakdown\_R3\\
NavierStokes.Comparator.navier\_stokes\_breakdown\_periodic.
\end{quote}
Their types quantify over every positive viscosity and retain the exact
initial-data, forcing, regularity, energy, and periodicity predicates.
The reference challenge modules intentionally contain unproved challenge
statements. They are kept separate from the solution imports. A source
scan that merely counts every occurrence of \texttt{sorry} in the whole
repository would therefore misdescribe the solution's dependency boundary.
Conversely, a source scan finding no such token cannot replace kernel
checking or comparison of the definitions with the intended mathematics.

The inspection found 816 files under the Navier--Stokes source directory
and 609 local modules in the final solution's import closure. Selected
modules were read in detail; all 609 were not manually reviewed. The
upstream metadata reports only the standard dependencies
\texttt{propext}, \texttt{Classical.choice}, and \texttt{Quot.sound} for
the final results. The accompanying verification ledger distinguishes
that upstream report from any fresh build and independent Comparator
execution performed in this extension \citep{NSLean2026}.

In this review the exact compiler distribution was downloaded and its
published SHA-256 checksum matched. Dependency cache retrieval then failed
with HTTP 429 responses after partial decompression. A direct attempt to
elaborate the final module stopped at missing local import artifacts.
Consequently this pass did not complete a fresh build of the final upstream
theorems, and did not run the independent Comparator certificate check.
The source audit found no challenge import and no unproved-placeholder token
in the 609-module local solution closure, after removing comments and string
literals. That lexical result has the narrower meaning stated here; it is
not a substitute for the uncompleted kernel verification.

The extension's new contribution is an explicit mathematical organization:
quadratic residual calculus; exact information-loss identities; norm
scaling that separates energy from regularity; a forbidden-word estimate
that explains a common analytic domain; all-order realization with
correct quantifiers; potential-based localization; and a quantitative
residual comparison for learned fields. Chapter~\ref{ch:fluid-structures}
then develops the arithmetic, exterior-power, and filtered-algebra
connections, while Chapter~\ref{ch:quantum-fluid} treats quantum
representations and learning.

These are rigorous bridges because each has a specified map or estimate.
They do not identify a finite surface-code complex with the continuum
Navier--Stokes equation, or turn a learned function into an existence
theorem without a stability certificate. They also preserve the released
theorem's force. Removing that force, imposing a different physical
constitutive law, or deriving the same singular flow from a microscopic
quantum state would be a further theorem with additional analytic content.
